% Transcription — fonds Alexandre Grothendieck, shelfmark 102, batch 1 (pages 1-13).
% Established from the facsimile archives/batches/102/batch-01.pdf (a mirror of
% https://grothendieck.umontpellier.fr/102.pdf), per the skill
% transcribe-grothendieck. The folder is thirteen pages and this is its only
% batch. Pages 7 (blank) and 9 (an unrelated printed verso, upside down,
% exercises on valuations with no hand on it) are skipped, the gap in the page
% numbers being the record. Page 1 is an orange cover inscribed « Recollements »
% in his hand; following hand.md §5 it gets no \page{} and its word heads the
% first section.
%
% The folder holds two things. Pages 2 to 6 are a typescript headed « Pour
% EGA VI », « Recollements d'objets d'un site »: a site S (standard, i.e.
% topology coarser than the canonical one — his interlinear insertion) with a
% class M of arrows satisfying a) stability and quarrability, b) identities and
% composition, c) locality on the target, d) a saturation condition (locality on
% the source, which he rewrites by hand); a Proposition extending M uniquely to a
% class M~ of arrows of the topos S~ (conditions a~ to d~ plus e~, f~) with a
% typed proof by cases on the representability of the diagonal; the Definition
% of objects of S~ « obtained by gluing objects of S along morphisms of M » and
% the full subcategory S~_M; a stability Proposition (i)-(v) and a Corollary
% (S~_M is the smallest class containing Ob S stable under sums and quotients by
% equivalence relations with pr_1 in M~, and is unchanged under replacing S by
% S' ⊂ S~_M); Examples (open immersions give schemes; étale morphisms give
% Artin's algebraic spaces); and a Proposition under a « condition
% draconienne » identifying the localised topos S~_M/X with the topos of
% S(X,M). Following folders 48 and 55, the statements and his interventions
% are given and the typed proofs are summarised in notes rather than
% recomposed. Page 5 is the verso of page 4 and carries only two typed lines,
% struck through, which continue the struck « Proposition 1) » at the foot of
% page 3; they are given as struck.
%
% Pages 8, 10, 11, 12 and 13 are manuscript notes, page 8 in black ink, the
% others in blue. Page 8 is « Questions pour Artin »: effectivity of descent
% data for proper schemes over a henselian local base, fpqc descent of
% « variétés », representability of proper algebraic spaces over a complete
% local base, and fpqc descent of small étale sheaves, with two marginal notes
% written slantwise. Page 10 is a numbered plan (1 to 8, with 7 moved by an
% arrow to follow 5) for a general study of functors on (Sch): representable
% morphisms, quasi-compact and locally finitely presented morphisms, formally
% smooth/unramified/étale, monomorphisms and immersions, separation, surjective
% and radicial morphisms, properness, connectedness, and the principle of
% verification for morphisms of sheaves. Page 11 is a fragment of a proof, taken
% up mid-sentence, reducing to a reduced fibre by a radicial extension and
% using M = Hilb_{X/Y} and a set-theoretic section φ; a boxed block is struck.
% Page 12 states that an fpqc sheaf monomorphism, infinitesimally étale,
% between functors commuting with filtered limits of rings and adic limits of
% artinian rings is an open immersion. Page 13 is a sheet of boxed keywords
% (a map of the planned classes of morphisms and conditions), transcribed as a
% list since its boxes do not carry arrows between them beyond one bracket.
%
% The manuscript hand is fast throughout, pages 8 and 11 worst: statements
% carry, connective words often do not, and what the leaves will not support is
% left \ill{}. The typescript is clean; the typist's own x-outs are not
% reproduced, only the strikings and insertions in his hand.
%
% Notation fixed once here. The typescript's tilde over a letter (S~, M~,
% S~_M, and the condition labels a~ ... f~, which he re-marks by hand) is set
% \widetilde{} / \tilde{}; the typed « x » of fibre products is \times; « sss »
% (si et seulement si) stays as typed; « ∐ » is \coprod. In the manuscript,
% « fpqc » is written with the letters run together; « loc. noeth. »,
% « loc. de prés. finie », « rel^t », « qu.-cpt » stay abbreviated as written.
%
% Readings flagged and not settled: the words between « S' → S » and « plat »
% in 1°) of page 8 and the whole of its left-margin notes; « séparable » on page
% 11 line 2; « compléments » (or « conditions ») in the bracket note of page 10;
% « induits » in item 7 of page 10; and « relatif » boxed on page 13.
%
% Pass: Opus 5 (claude-opus-5), 2026-09-14 — first pass, unchecked against the pages by a human.
\documentclass[11pt,a4paper]{article}
% Le préambule commun à toutes les transcriptions du fonds Grothendieck.
%
% Il tient en une idée : **l'incertitude se compose, elle ne se résout pas.**
% Une transcription qui lisse un mot douteux a détruit la seule chose pour
% laquelle on la faisait — un lecteur doit pouvoir savoir, à chaque endroit,
% ce qui était sur le feuillet et ce qui a été ajouté. Les six macros qui
% suivent sont donc l'appareil critique, et non de la décoration.
%
% Les mêmes commandes sont reconnues par `scripts/render.mjs`, qui produit la
% vue de lecture du site. Une macro ajoutée ici doit l'être aussi là-bas,
% faute de quoi le rendu échouera bruyamment — ce qui est voulu : un rendu
% silencieusement faux est pire qu'un rendu absent.

\ProvidesPackage{grothendieck}[2026/08/08 Transcriptions du fonds Grothendieck]

\usepackage{amsmath,amssymb,amsthm}
\usepackage{mathrsfs}
\usepackage{stmaryrd}
% Les diagrammes commutatifs sont partout dans ce fonds : c'est la langue
% même de la géométrie algébrique de Grothendieck, et une transcription qui
% les rendrait en prose ne transcrirait rien.
\usepackage{tikz-cd}
% Français par défaut — c'est la langue des feuillets — mais l'anglais est
% chargé aussi : la traduction et le résumé sont des documents anglais, et la
% typographie française (espaces avant les deux-points) y serait une faute.
% Ces documents basculent par \selectlanguage{english} en tête de corps.
\usepackage[english,french]{babel}
\usepackage{fontspec}
\usepackage{geometry}
\geometry{a4paper,margin=25mm,marginparwidth=28mm}
\usepackage{marginnote}
\usepackage{xcolor}
% ulem, not soul. `soul`'s \st and \ul are built on a character-by-character
% scan that XeLaTeX's font machinery breaks: the document fails with an
% undefined control sequence rather than degrading. ulem does the same two jobs
% and is Unicode-engine safe. `normalem` stops it hijacking \emph, which the
% transcriptions use for Grothendieck's own emphasis.
\usepackage[normalem]{ulem}
\usepackage{hyperref}
\hypersetup{colorlinks=true,linkcolor=black,urlcolor=black,pdfborder={0 0 0}}

% --- Métadonnées du lot -----------------------------------------------------
% Reprises telles quelles par le script de rendu, qui les affiche en tête de
% la vue de lecture. Elles fixent ce que le fichier prétend couvrir : sans
% elles, une transcription flotte sans point d'attache dans un fonds de seize
% mille feuillets.

\newcommand{\folder}[1]{\gdef\grothFolder{#1}}
\newcommand{\batch}[1]{\gdef\grothBatch{#1}}
\newcommand{\leaves}[2]{\gdef\grothFirst{#1}\gdef\grothLast{#2}}
\newcommand{\foldertitle}[1]{\gdef\grothFoldertitle{#1}}
\gdef\grothFolder{?}\gdef\grothBatch{?}\gdef\grothFirst{?}\gdef\grothLast{?}\gdef\grothFoldertitle{}

% --- Le résumé --------------------------------------------------------------
%
% Il ouvre la lecture modernisée, sous le titre « Résumé », et il est composé
% en petit. Ce n'est pas un ornement : c'est le seul passage du document qui
% ne soit pas des mathématiques, et un lecteur doit pouvoir le reconnaître
% avant de l'avoir lu — puis le sauter s'il connaît déjà le sujet, ou s'y
% arrêter s'il ne le connaît pas. Le filet qui le suit marque la reprise du
% corps.

\newenvironment{resume}
  {\small\setlength{\parskip}{0.45em}}
  {\par\medskip\hrule\medskip}

% --- Mots-clés ---------------------------------------------------------------
%
% En anglais, en clôture du résumé de la lecture modernisée : le vocabulaire
% moderne sous lequel chercher ce que les feuillets construisent. Le manifeste
% du site (`npm run manifest`) les extrait pour étiqueter la cote entière —
% cette ligne est la source unique des tags, il n'y a pas de fichier à côté.
\newcommand{\keywords}[1]{\par\smallskip\noindent
  {\footnotesize\textsc{Keywords} --- \textit{#1}}\par}

% --- L'appareil critique ----------------------------------------------------

% Le numéro de feuillet, en marge. C'est l'ancre qui permet de régler un
% désaccord sur une formule en regardant *un* feuillet plutôt qu'une chemise
% de six cents. Le site s'en sert aussi pour tourner les pages du fac-similé
% quand on fait défiler la transcription.
\newcommand{\leaf}[1]{%
  \marginnote{\footnotesize\color{gray!70}\textsf{#1}}%
  \label{leaf:#1}\ignorespaces}

% Illisible : on ne devine pas. Un mot inventé qui se lit comme les autres est
% le pire résultat possible.
\newcommand{\ill}{\textcolor{red!60!black}{[\,\ldots\,]}}

% Lecture douteuse : le mot est donné, et signalé comme incertain.
\newcommand{\uncertain}[1]{\uline{#1}}

% Ajout de l'éditeur — une lettre manquante, un mot sous-entendu.
\newcommand{\add}[1]{\textcolor{blue!50!black}{[#1]}}

% Biffé par Grothendieck lui-même. Ce qu'il a rayé fait partie du document :
% une rature dit souvent où la pensée a changé de direction.
\newcommand{\struck}[1]{\sout{#1}}

% Note du transcripteur, sur l'état matériel du feuillet.
\newcommand{\note}[1]{\par\smallskip{\small\color{gray!60!black}\textit{[#1]}}\par\smallskip}

% Note marginale de Grothendieck, distincte du corps du texte.
\newcommand{\marginal}[1]{\par\smallskip\noindent\hspace{1em}%
  {\small\color{gray!60!black}#1}\par\smallskip}

% --- Titre ------------------------------------------------------------------

\AtBeginDocument{%
  \begin{center}
    {\large\bfseries Fonds Alexandre Grothendieck --- cote n\textsuperscript{o}~\grothFolder}\\[2pt]
    {\itshape\grothFoldertitle}\\[4pt]
    {\small feuillets \grothFirst--\grothLast\ (lot \grothBatch)}\\[2pt]
    {\footnotesize\color{gray!60!black}Transcription établie d'après le fac-similé numérisé
     par l'Université de Montpellier.\\
     Les lectures incertaines sont soulignées, les illisibles notées [\,\ldots\,],
     les ajouts entre crochets.}
  \end{center}
  \vspace{1em}}

\endinput


\folder{102}
\batch{1}
\pages{1}{13}
\dating{[après 1967]}
\watermark{Édition de démonstration}
\foldertitle{Etude générale des foncteurs (Sch)° → (Ens) : tapuscrit annoté (s.d.), notes manuscrites (s.d.).}

\begin{document}

\section*{Recollements}

\note{le mot est de sa main, seul, en haut à droite de la chemise orange qui
forme la page 1 ; il est repris ici comme titre de section}

\page{2}

\note{feuillet de tapuscrit, repris à la main à l'encre noire. Comme pour les
tapuscrits des dossiers 48 et 55, on donne les énoncés et ses interventions ;
les démonstrations tapées sont résumées}

Pour EGA VI

\subsection*{Recollements d'objets d'un site.}

Soit $S$ un site\supplied{ standard (i.e.\ à topologie moins fine que la
canonique)}, muni d'une partie $M$ de l'ensemble des flèches, satisfaisant les
conditions habituelles :

\begin{itemize}
\item[a)] les flèches de $M$ sont quarrables et stables par changement de base,

\item[b)] $M$ contient les flèches identiques et est stable par composition,

\item[c)] une flèche $u : X \to Y$ appartient à $M$ si elle y appartient
localement sur $Y$,
\end{itemize}

et en plus une condition de saturation :

\begin{itemize}
\item[d)] Soit $u_i : X_i \to X$ une \add{famille couvrante de}
\struck{$Y$-}morphismes dans \uncertain{$S$}, \struck{supposons $X_i \to Y$
dans} \add{et supposons les $uu_i : X_i \to Y$ dans $M$, (i.e., par a) et c),
les} $X_i \times_X X_j \xrightarrow{\mathrm{pr}_1} X_i$ dans $M$), alors
$u$\add{$: X \to Y$} est dans $M$.
\end{itemize}

\note{la condition d) est entièrement réécrite à la main : « famille couvrante
de » au-dessus de la ligne, le « $Y$- » de « $Y$-morphismes » biffé, « supposons
$X_i \to Y$ dans » tapé et biffé, remplacé en interligne par « et supposons les
$uu_i : X_i \to Y$ dans $M$, », et au-dessus encore « (i.e.\ par a) et c), les »
qui fait de la ligne tapée $X_i \times_X X_j \to X_i$ dans $M$ une parenthèse
explicative ; le $Y$ en indice du produit fibré tapé est corrigé en $X$.
« $: X \to Y$ » est ajouté sous « $u$ ». Le « $S$ » après « dans » est surchargé
d'un trait vertical et pourrait avoir été changé}

\marginal{essayer de ne pas le mettre ici, en saturant plus tard en
$\widetilde{M}$, et formulant $\tilde{e}$) de façon un peu différente : $u$ est
dans $M$ sss on peut trouver des $u_i : X_i \to X$ dans $M$ avec $uu_i$ dans
$M$}\note{note écrite de biais dans la marge gauche, en regard de d) et du
début de la proposition, et reliée à d) par une accolade}

\textbf{Proposition.} Il existe une partie $\widetilde{M}$ de l'ensemble des
flèches du topos $\widetilde{S}$, déterminée de façon unique par les propriétés
suivantes : $\widetilde{M}$ satisfait aux conditions \add{$\tilde{a}$,
$\tilde{b}$, $\tilde{c}$, $\tilde{d}$} analogues à a) à d), et de plus aux
conditions suivantes :

\begin{itemize}
\item[$\tilde{e}$)] Si $u : X \to Y$ est une flèche de $S$, alors $u$ est dans
$M$ sss $u$ est dans $\widetilde{M}$.

\item[$\tilde{f}$)] Si $u : F \to X$ est une flèche dans $\widetilde{M}$, avec
$X \in S$\note{le $S$ de « $X \in S$ » est surchargé à la main et se lit mal}, alors il existe une famille couvrante $u_i : X_i \to F$
de flèches dans $\widetilde{M}$, avec $X_i \in S$. \struck{\ill{}}
\end{itemize}

\note{« $\tilde{a}, \tilde{b}, \tilde{c}, \tilde{d}$ » est écrit à la main
au-dessus de « analogues à a) à d) ». À la fin de $\tilde{f}$) quelques mots
tapés ont été biffés à la main ; on y devine « $uu_i$ … $\widetilde{M}$ »}

Démontrons l'unicité, et montrons que l'on peut décider sans ambiguïté si une
flèche $u : F \to G$ de $\widetilde{S}$ est dans $\widetilde{M}$. Pour ceci,
distinguons deux cas :

1º) Cas où le morphisme diagonal $F \to F \times_G F$ est représentable.

\note{la démonstration tapée qui suit, et qui se poursuit sur la page 3, est
résumée : on couvre $F$ par des $X_i \in S$, et par $\tilde{c}$) on se ramène
au cas où $G = X \in S$ ; on montre alors que $u \in \widetilde{M}$ sss il existe
une famille couvrante $u_i : X_i \to F$, $X_i \in S$, avec $uu_i \in M$ et les
$X_i \times_F X_j \xrightarrow{\mathrm{pr}_1} X_i$ dans $M$}

\page{3}

\note{suite du tapuscrit. Ses interventions sur cette page sont : les
étiquettes des conditions remarquées d'un tilde à la main
($\tilde{b}$, $\tilde{c}$, $\tilde{e}$, $\tilde{f}$, $\tilde{a}$, $\tilde{d}$) ;
dans la parenthèse « (NB cela a un sens, car l'hypothèse faite sur $u$ implique
que $X_i \times_F X_j$ est représentable, comme image inverse de la diagonale de
$F$ sur $X$) », l'insertion « dans $X_i \times_X X_j$ » au-dessus de « de la
diagonale » ; et, dans le cas 2º), « morphisme diagonal » écrit au-dessus d'un
mot tapé biffé. Le cas 2º) (cas général) se ramène au cas 1º) en observant que
$Z_{ij} = X_i \times_F X_j$ est un sous-faisceau du faisceau représentable
$Z'_{ij} = X_i \times_X X_j$, image inverse du morphisme diagonal, donc
représentable. Le tapuscrit laisse au lecteur la vérification que
$\widetilde{M}$ satisfait à $\tilde{a}$) à $\tilde{f}$)}

\textbf{Définition.} On dit qu'un objet $F$ de $\widetilde{S}$ peut s'obtenir
\emph{par recollement d'objets de $S$ à l'aide de morphismes dans $M$}, si on
peut trouver une famille couvrante $u_i : X_i \to F$, avec $X_i \in S$ et
$u_i \in \widetilde{M}$. On désignera par $\widetilde{S}_M$ la sous-catégorie
pleine de $\widetilde{S}$ ayant comme objets ces $F$.

\note{les mots « par recollement d'objets de $S$ à l'aide de morphismes dans $M$ »
sont soulignés à la main}

On a alors la propriété de stabilité suivante :

\struck{\textbf{Proposition} 1) Soit $F \in \mathrm{Ob}\,\widetilde{S}$, et
$F_i \to F$ une famille couvrante}

\page{4}

\textbf{Proposition.} (i) Tout faisceau représentable $F \in S$ est dans
$\widetilde{S}_M$.

(ii) Soit un morphisme $u : F \to G$ dans $\widetilde{M}$; si
$G \in \mathrm{Ob}\,\widetilde{S}_M$, alors $F \in \mathrm{Ob}\,\widetilde{S}_M$.
La réciproque est vraie si $u$ est couvrant.

(iii) Soit $u : F \to G$ un morphisme couvrant, tel que
$F \times_G F \xrightarrow{\mathrm{pr}_1} F$ soit dans $\widetilde{M}$, alors
$G$ est dans $\widetilde{S}_M$ sss $F$ l'est.

(iv) Une somme $\coprod F_i$ de faisceaux est dans $\widetilde{S}_M$ sss les
$F_i$ le sont.

(v) Supposons que $S$ soit stable par produit de deux objets (resp.\ par
produit fibré de deux objets sur un troisième, resp.\ par objet final, resp.\
par $\varprojlim$ finies) alors il en est de même de $\widetilde{S}_M$.

\note{« $\mathrm{pr}_1$ » dans (iii) est écrit à la main au-dessus de la
flèche ; le numéro « (iv) » est une surcharge d'un « (iii) » tapé ; l'énoncé (v)
est retapé en interligne sur une première rédaction biffée à la machine
(« Soit I une catégorie finie, telle que … »)}

(i) est trivial, ainsi que la réciproque dans (ii). La partie directe est
facile ….

\note{la démonstration tapée est résumée : (iii) résulte de (ii) ; pour (iv),
si $X \to F = \coprod F_i$ est dans $\widetilde{M}$ avec $X$ représentable, les
$X_i = X \times_F F_i$ se couvrent par des $X_{ij} \in S$ par $\tilde{f}$), et
les $X_{ij}$ couvrent les $F_i$. Pour (v), le tapuscrit renvoie au rédacteur}

On laisse aussi (v) au rédacteur (attention, il est douteux que ça marche pareil
pour un noyau, car il ne sera pas vrai nécessairement qu'un morphisme de
doublets dans $\widetilde{S}$ \add{ou même de $S$} induisant des morphismes
$\in \widetilde{M}$ sur les sommets induise un morphisme $\in M$ sur les noyaux
des double flèches).

\textbf{Corollaire.} $\mathrm{Ob}\,\widetilde{S}_M$ est le plus petit
sous-ensemble de $\mathrm{Ob}\,\widetilde{S}$ contenant $\mathrm{Ob}\,S$ et
stable par \add{sommes et par} passage au quotient par une relation d'équivalence
à $\mathrm{pr}_1$ qui est dans $\widetilde{M}$. Par suite, $\widetilde{S}_M$ ne
change pas si on remplace $S$ par une sous-catégorie $S'$ de $\widetilde{S}$
contenue dans $\widetilde{S}_M$, et $M$ par
$M' = \widetilde{M} \cap \mathrm{Fl}\,S'$.

\note{le tapuscrit porte « le plus petit sous-ensemble de $\mathrm{Ob}\,S$ »,
sans tilde ; le sens veut $\widetilde{S}$ et le tilde n'est pas tracé à cet
endroit}

\page{5}

\struck{telle que les $F_i \times_F F_j \to F_i$ soient dans $\widetilde{M}$.
Pour que l'on ait $F \in \mathrm{Ob}\,\widetilde{S}_M$ il f\supplied{au}t
\add{et il suffit} que pour tout $i$, on ait $F_i$}

\note{ces deux lignes tapées, au haut du verso de la page 4, continuent la
proposition biffée au bas de la page 3, et sont barrées de plusieurs traits
obliques ; la version retenue est la proposition de la page 4. Le « il fet s. »
tapé est lu « il faut et il suffit »}

\page{6}

\textbf{Exemples :} Si $S$ est la catégorie des schémas affines, avec une
topologie plus fine que Zar, $M$ l'ensemble des immersions ouvertes (ou, en
l'absence de d), les immersions de la forme
$\mathrm{Spec}(A_f) \to \mathrm{Spec}(A)$), alors la construction donne une
catégorie $\widetilde{S}_M$ équivalente à celle des schémas. Si on prend une
topologie plus fine que la topologie étale, et pour $M$ \add{l'ensemble} des
morphismes étales (qui satisfont la condition d)), alors on trouve les
« espaces algébriques » de Artin, en « recollant des spectres d'anneaux par des
morphismes étales ». Le dernier corollaire nous dit qu'on aurait trouvé le même
résultat en partant de la catégorie $S'$ des schémas, au lieu de celle $S$ des
schémas affines. \add{Dans ce cas, les morphismes de $\widetilde{M}$ se disent
\emph{étales} [la top.\ sur $S$ étant la top.\ étale).}

\note{la dernière phrase est ajoutée à la main, à l'encre bleue, à la suite de
la ligne tapée}

\textbf{Proposition.} Supposons que $M$ satisfasse la condition draconienne
suivante : soient $u : X \to Y$ et $v : Y \to Z$ des flèches composables dans
$S$, alors $v \in M$ et $vu \in M$ implique $u \in M$. Pour un objet $X$ de $S$,
soit $S(X,M)$ la sous-catégorie de $S_{/X}$ formée des $X' \to X$ qui sont dans
$M$ (donc la condition précédente dit que tous les morphismes de $S(X,M)$ sont
dans $M$); munissons-la de la topologie induite par $S$, et considérons le
topos $S(X,M)^{\sim}$. Alors \struck{le foncteur
$(\widetilde{S}_M)_{/X} \to S(X,M)^{\sim}$ composé de
$(\widetilde{S}_M)_{/X} \xrightarrow{\text{incl}} \widetilde{S}_{/X}
\xrightarrow{\text{can.}\ \approx} (S_{/X})^{\sim}
\xrightarrow{\text{restr}} S(M,X)^{\sim}$ est une équivalence de catégories.}
la sous-catégorie pleine de $\widetilde{S}_{/X}$ formée des $F \to X$ qui sont
dans $\widetilde{M}$ \add{[d'où résulte $F \in \mathrm{Ob}\,\widetilde{S}_M$ !]},
est un \emph{topos}, équivalent par le foncteur restriction à $S(X,M)$ au topos
$S(X,M)^{\sim}$.

\note{la première rédaction de la conclusion, tapée, est rayée à la main d'un
trait horizontal et d'un zigzag, et encadrée ; dans le composé biffé, le
dernier terme est tapé « $S(M,X)$ ». L'insertion entre crochets est de sa main à
l'encre bleue, au-dessus de « dans $\widetilde{M}$ ». Dans la dernière ligne,
« $S(X,M)$ » est une surcharge d'une frappe erronée}

\section*{Questions pour Artin}

\page{8}

\note{le titre, souligné, est de sa main en haut de la page}

1º) $S$ \struck{strictement} local \struck{noeth.} \emph{hens.} noeth.,
$S' \to S$ \struck{\ill{}} \uncertain{plat.} \emph{surjectif}, $S'$ noeth.,
$X'/S'$ \struck{loc.\ de t.f.} \add{propre} avec donnée de descente rel.\ à
$S'/S$. Supposons que la donnée de descente sur $X'_s$ soit effective ($s$ pt.\
fermé de $S$). Alors la donnée de descente sur $X'$ est effective ?

\note{« hens. » est doublement souligné. « propre » est écrit au-dessus de
« loc.\ de t.f. », biffé}

Cas où $S' = \widehat{S}$ ?

2º) Peut-on faire la descente fpqc de \emph{variétés} \add{ou relation loc.\
de prés.\ finie,} p.\ ex.\ de parties d'un complété d'un anneau local ? [i.e.\
la propriété, pour un faisceau fpqc loc.\ de prés.\ finie, d'être une variété
est-elle loc.\ fpqc sur $S$ ?]

\note{« ou relation loc.\ de prés.\ finie, » est une insertion au-dessus de la
ligne, rattachée par un trait à « variétés » ; « relation » est une lecture
possible, « relativement » en serait une autre. La parenthèse entre crochets
est écrite en deux lignes serrées, la seconde sous la première}

\marginal{Le Pb \ill{} \uncertain{le comm.} aux $\varprojlim$
(\uncertain{adiques artiniennes})}\note{note écrite de biais dans la marge
gauche, en regard de 2º) ; « $\varprojlim$ » est un « lim » souligné d'une
flèche vers la gauche}

3º) $S$ comme dans 1º, $X$ une \emph{variété} \struck{sur} \struck{loc.\ de
t.f.} \add{propre} sur $S$, \add{et $X_s$ représentable,} \add{alors $X$
est-il} \struck{si $X$ est représentable loc.\ fpqc sur $S$,} représentable ?
Évidemment, si la réponse à 2º \struck{est} était affirmative, les questions 1º
et 3º seraient équivalentes…

\note{la rédaction de 3º) est reprise en interligne : « et $X_s$
représentable, » sous la première ligne, « alors $X$ est-il » encadré au-dessus
de la seconde ; la lecture du passage biffé entre les deux reste douteuse}

\marginal{Exemples de $X$ \struck{variétés} propre sur $S$ local complet, $X_s$
repr.\ mais $X$ non repr.\ ?; Exemples de $X$ variétés \uncertain{propres} sur
$S$ noeth., \ill{} \uncertain{repr.}, qui ne soit repr.\ ?}\note{note écrite de
biais dans la marge gauche en regard de 3º) et 4º), en deux paragraphes
commençant chacun par « Exemples » ; plusieurs mots en sont incertains}

4º) Descente fpqc des \add{(« petits »)} faisceaux étales.

Lié à la question des conditions sur les foncteurs
$(\mathrm{Sch})^{\circ}_{/S} \to (\mathrm{Ens})$ qui proviennent de petits
faisceaux étales sur $S$. Une réponse affirmative à 2º (i.e.\ \uncertain{pour
des} « variétés » dans le sens \uncertain{généralisé}, sans exiger que
\add{la} diagonale soit une \add{immersion}) en impliquerait une à 4º.

\note{« 4º » est entouré, et un trait le relie à la note marginale}

\section*{Plan}

\note{titre de l'édition : la page 10 ne porte pas de titre}

\page{10}

\begin{enumerate}
\item[1.] Morphismes représentables. Morphismes $P$-représentables.

Abus de langage : imm., imm.\ ouverte, imm.\ fermée, lisse, net, étale,
(sous-entendu : repr.\ par …). Critère de représentabilité par
« recollement des morceaux ». (SGAD XI 3.5.)

\item[2.] [Faisceaux qu.-cpt, quasi-sép.\ Morphismes qu.-cpt, qu.-sép.\ …\
des faisceaux]. Morphismes \add{loc.}\ de prés.\ finie, de prés.\ finie.

\marginal{[ ] = petits \uncertain{compléments}}

\item[3.] Morphismes formel\supplied{lemen}t lisses, formel\supplied{lemen}t nets,
formel\supplied{lemen}t étales, et variantes infinitésimales. Critères standard…

\item[4.] Critères pour qu'un morphisme soit un monomorphisme, un bimorphisme,
une immersion, une immersion ouverte.

\item[5.] Conditions de séparation : \emph{morph.\ diag.\ repr.}\ par imm.,
par imm.\ qu.-cpte, [par imm.\ de prés.\ finie] ; morphisme diagonal repr.\ par
imm.\ fermée ; Relations avec critères \struck{valutifs} \uncertain{valuatifs}.

\item[7.] Espaces \uncertain{induits} et foncteurs \struck{\ill{}} sur
$(\mathrm{Sch})_{/\ill{}}$. Morphismes surjectifs, \struck{radiciels} radiciels
(morphisme diagonal est surjectif). [Une immersion ouverte surjective est un
isom].

\item[6.] Conditions de propreté.

\item[8.] Conditions de connexions.
\end{enumerate}

\note{une flèche tracée dans la marge ramène l'item 7 entre les items 5 et 6 ;
l'ordre de la page est gardé. Les trois lignes de 5, ici séparées par des
points-virgules, sont réunies par une accolade ; « par imm.\ de prés.\ finie » est encadré et relié à « par imm.\
qu.-cpte ». Autour de « critères valutifs » la page est surchargée : un mot
biffé, et « Morphismes », d'une encre plus sombre, écrit par-dessus la fin de la
ligne, et l'on n'est pas sûr de la frontière entre 5 et 7}

Principe des vérifications pour les morphismes de faisceaux sur
$(\mathrm{Sch})$ : a) Conditions locales b) Si $f : F \to G$ est
représentable, cela doit coller avec la définition déjà acceptée pour morphismes
de schémas [$\forall\, X \to G$, le morphisme de schémas $X \times_G F \to X$ a
la propriété].

\section*{Fragment d'une démonstration}

\note{titre de l'édition ; la page 11 reprend au milieu d'une phrase, dont le
début ne figure pas dans le dossier}

\page{11}

dans une extension radicielle convenable de \uncertain{$k(y)$}, on peut supposer
que $(X_y)_{\mathrm{red}}$ est \uncertain{séparable} sur \add{$k(y) =$}
\uncertain{$k(y)$}, et remplaçant $X$ par $X_{\mathrm{red}}$, on \ill{} \ill{}
les conditions du th., et on gagne.

La deuxième assertion se déduit de la troisième par réduction au cas
noethérien. Pour la troisième,

\struck{distinguons les deux cas. \ill{} \ill{} \ill{} $Y$ \ill{} \ill{} \ill{}
$Y$ intègre, \ill{} \ill{} $Y$. Faisant une extension radicielle finie $Y'$ de
$Y$ et normalisant \ill{} $Y$ réduit. Dans le premier cas.}

\note{bloc encadré sur la gauche et biffé de plusieurs traits horizontaux et
obliques ; seuls quelques mots se lisent}

(On \uncertain{peut supposer} $Y$ réduit.) \struck{\ill{} \ill{}
$M = \underline{\mathrm{Hilb}}$}

\struck{\ill{} \ill{}} \uncertain{Notons} $M = \underline{\mathrm{Hilb}}_{X/Y}$,
\ill{} alors une section \emph{ensembliste} $\varphi$ de $M$ sur $Y$,
\struck{telle que} définie par : $\forall\, y \in Y$ le point
\add{$\varphi(y)$} de $M_y$ \uncertain{défini par}
$(X_y)_{\mathrm{red}} \subset X_y$. \struck{Soit $D'$} Soit
$Y' = \varphi(Y)$. Je dis que $Y'$ est une partie \struck{[\ill{} \ill{}]}
\ill{} \ill{} $Y'$ de la \uncertain{fermée} de $M$; \struck{\ill{}}
\uncertain{admettons} \ill{} [NB. $Y' \to Y$ est \ill{} \ill{}, \ill{} fini…]
\struck{\ill{} \ill{} admis}, \ill{} aura un morphisme fini \ill{}
\uncertain{fermé}, on \struck{$\varphi$} \ill{} aura un $Y$-morphisme
$Y' \to M$, correspondant à un sous-schéma \add{\uncertain{fermé}}
\struck{$X$} $Z$ de $X'$, plat de prés.\ finie sur $Y$, dont les fibres sont
\ill{} $(X_y)_{\mathrm{red}} \otimes_{k(y)} k(y')$, qui sont réduites. Par
suite, $Z$ est \uncertain{réduit} sur $Y'$ \add{réduit}, donc
\uncertain{évident}, et ayant même ensemble sous-jacent que $X'$, il est égal
à $X'_{\mathrm{red}}$].

\note{les lignes du milieu, de « Je dis que » à « un $Y$-morphisme », sont
surchargées d'insertions interlinéaires et de biffures ; on n'y a gardé que ce
qui se lit, et l'ordre des mots rendus n'est pas assuré. La parenthèse ouverte
par « [NB. » se ferme au bas de la page, après « $X'_{\mathrm{red}}$ ». « sur
$Y$ » dans « plat de prés.\ finie sur $Y$ » est lu tel qu'écrit, bien que le
sens attende $Y'$}

\section*{Immersions ouvertes}

\note{titre de l'édition}

\page{12}

\begin{tikzcd}
G \arrow[r, "u"] & F \arrow[d] \\
 & S
\end{tikzcd}

$S$ loc.\ noeth.

\struck{Si \ill{}} \add{Supposons} :

(i) $F$ et $G$ sont des faisceaux fpqc, commutant aux $\varinjlim$ d'anneaux,
aux $\varprojlim$ \uncertain{adiques} d'anneaux artiniens,

(ii) $u$ est un monomorphisme, et est infinitésimalement étale.

\struck{(iii) $u$ \ill{} \ill{} épimorphisme}

Alors $G \to F$ est une immersion ouverte.

\note{« $\varinjlim$ » et « $\varprojlim$ » sont des « lim » soulignés d'une
flèche vers la droite et vers la gauche. Les deux énoncés suivants sont séparés
du reste par des traits horizontaux}

Une Immersion ouverte surjective de foncteurs est un isomorphisme

Caractérisation des morphismes étales de foncteurs…,

\page{13}

\note{feuille de mots-clés encadrés, qui semble dresser la carte des notions à
traiter ; les cadres ne sont pas reliés par des flèches, sauf une accolade. On
les donne en liste, dans l'ordre de lecture}

$X \to Y$ une imm.\ ouverte

$X \to Y$

morphisme \struck{affine} formel\supplied{lemen}t lisse, non ram., étale ---
rel\supplied{ativemen}t lisse, rel\supplied{ativemen}t non ram., rel\supplied{ativemen}t étale

\textbf{Mon} : immersion; \struck{immersion ouverte}; immersion \emph{fermée}
--- imm. ouverte

[séparé, —] [propre, …] [connexe, …]

\note{« formellement » est surchargé sur un « affine » biffé. « Mon » est dans
un petit cadre relié au cadre des immersions}

$F_i \to F$  \emph{imm.\ ouvert\supplied{e}}

\begin{tikzcd}
G \arrow[r] & F \\
G_T \arrow[r] \arrow[u] & F_T \arrow[u]
\end{tikzcd}

(A)  \struck{\uncertain{Ex.}}  faisceaux; \emph{loc.\ de
prés.\ finie} [\uncertain{relatif}]; adiques

\note{les trois derniers cadres (« faisceaux », « loc.\ de prés.\ finie
[relatif] », « adiques ») sont réunis par une accolade à laquelle aboutit un
mot biffé ; le « A » est entouré}

\end{document}
